{\actuality} В эпоху цифровой трансформации стремительное развитие технологий блокчейн
открывает новые горизонты для бизнеса и государственных учреждений. Современные
организации все чаще стремятся интегрировать распределенные реестры для повышения
прозрачности, безопасности и эффективности своих операций. Однако текущие решения
зачастую требуют значительных ресурсов на внедрение и сопровождение. Актуальность
настоящего исследования обусловлена необходимостью упрощения процесса создания и
управления блокчейн-сетями. Объединение возможностей Hyperledger Fabric и сети
Cosmos может предложить более гибкое и экономически эффективное решение. Это
особенно важно для организаций, стремящихся оптимизировать бизнес-процессы без
значительных затрат.

{\aim} данной работы является детальное изучение и разработка механизма, который
позволяет создавать блокчейн-сети на базе уже существующих валидаторов, аналогично
архитектуре, реализованной в Hyperledger Fabric, с последующей адаптацией для
экосистемы Cosmos/Tendermint.

Для достижения поставленной цели необходимо решить следующие {\tasks}:
\begin{enumerate}[beginpenalty=10000]
  \item Проанализировать эволюцию блокчейн-технологий и ключевые принципы их работы.
  \item Исследовать архитектуру и возможности платформы Cosmos и протокола Tendermint.
  \item Сравнить инфраструктуру Hyperledger Fabric, Ethereum и Polkadot с точки зрения
  построения частных и консорциумных сетей.
  \item Сформировать требования к автоматизированной инфраструктуре валидаторов и
  предложить схему их подбора на основе производительности.
  \item Разработать архитектуру восстановления смарт-аккаунтов с применением
  пороговой криптографии и социальной модели доверия.
  \item Интегрировать прототип и провести оценку масштабируемости и надежности.
  \item Сформулировать выводы и рекомендации для развертывания и развития решения.
\end{enumerate}

{\novelty}
\begin{enumerate}[beginpenalty=10000]
  \item Предложена архитектура автоматизированного формирования сети валидаторов
  в экосистеме Cosmos/Tendermint, ориентированная на повторное использование
  существующей инфраструктуры.
  \item Разработан подход к оценке валидаторов на основе производительности и
  надежности для динамического обновления состава сети.
  \item Сформулирована модель восстановления смарт-аккаунтов, совместимая с
  требованиями приватных блокчейнов и механизмами Tendermint.
\end{enumerate}

{\influence} Практическая значимость работы заключается в возможности упрощенного
развертывания частных и консорциумных сетей Cosmos с сокращением затрат на
сопровождение и повышением управляемости валидаторской инфраструктуры.

{\methods} В работе используются методы системного анализа, архитектурного
проектирования, сравнительного анализа платформ, моделирования и экспериментальной
оценки прототипа.

{\defpositions}
\begin{enumerate}[beginpenalty=10000]
  \item Архитектура автоматизированной сети валидаторов позволяет строить новые
  блокчейн-сети на базе существующих валидаторов без потери устойчивости.
  \item Схема отбора валидаторов на основе метрик производительности повышает
  надежность и масштабируемость сети.
  \item Механизм восстановления смарт-аккаунтов на основе пороговой криптографии
  обеспечивает баланс между доступностью и безопасностью.
\end{enumerate}

{\reliability} полученных результатов обеспечивается использованием формализованных
метрик, экспериментальной проверкой на прототипе и сопоставлением с результатами,
полученными в смежных исследованиях.

{\probation}
Основные результаты работы докладывались на профильных семинарах и конференциях
(будет уточнено по мере подготовки рукописи).

{\contribution} Автор выполнил анализ платформ, сформулировал архитектуру решения,
разработал модель оценки валидаторов и подготовил прототип для экспериментальной проверки.

\ifnumequal{\value{bibliosel}}{0}
{%%% Встроенная реализация с загрузкой файла через движок bibtex8. (При желании, внутри можно использовать обычные ссылки, наподобие `\cite{vakbib1,vakbib2}`).
    {\publications} Основные результаты по теме диссертации изложены
    в~XX~печатных изданиях,
    X из которых изданы в журналах, рекомендованных ВАК,
    X "--- в тезисах докладов.
}%
{%%% Реализация пакетом biblatex через движок biber
    \begin{refsection}[bl-author, bl-registered]
        % Это refsection=1.
        % Процитированные здесь работы:
        %  * подсчитываются, для автоматического составления фразы "Основные результаты ..."
        %  * попадают в авторскую библиографию, при usefootcite==0 и стиле `\insertbiblioauthor` или `\insertbiblioauthorgrouped`
        %  * нумеруются там в зависимости от порядка команд `\printbibliography` в этом разделе.
        %  * при использовании `\insertbiblioauthorgrouped`, порядок команд `\printbibliography` в нём должен быть тем же (см. biblio/biblatex.tex)
        %
        % Невидимый библиографический список для подсчёта количества публикаций:
        \phantom{\printbibliography[heading=nobibheading, section=1, env=countauthorvak,          keyword=biblioauthorvak]%
        \printbibliography[heading=nobibheading, section=1, env=countauthorwos,          keyword=biblioauthorwos]%
        \printbibliography[heading=nobibheading, section=1, env=countauthorscopus,       keyword=biblioauthorscopus]%
        \printbibliography[heading=nobibheading, section=1, env=countauthorconf,         keyword=biblioauthorconf]%
        \printbibliography[heading=nobibheading, section=1, env=countauthorother,        keyword=biblioauthorother]%
        \printbibliography[heading=nobibheading, section=1, env=countregistered,         keyword=biblioregistered]%
        \printbibliography[heading=nobibheading, section=1, env=countauthorpatent,       keyword=biblioauthorpatent]%
        \printbibliography[heading=nobibheading, section=1, env=countauthorprogram,      keyword=biblioauthorprogram]%
        \printbibliography[heading=nobibheading, section=1, env=countauthor,             keyword=biblioauthor]%
        \printbibliography[heading=nobibheading, section=1, env=countauthorvakscopuswos, filter=vakscopuswos]%
        \printbibliography[heading=nobibheading, section=1, env=countauthorscopuswos,    filter=scopuswos]}%
        %
        \nocite{*}%
        %
        {\publications} Основные результаты по теме диссертации изложены в~\arabic{citeauthor}~печатных изданиях,
        \arabic{citeauthorvak} из которых изданы в журналах, рекомендованных ВАК%
        \ifnum \value{citeauthorscopuswos}>0%
            , \arabic{citeauthorscopuswos} "--- в~периодических научных журналах, индексируемых Web of~Science и Scopus%
        \fi%
        \ifnum \value{citeauthorconf}>0%
            , \arabic{citeauthorconf} "--- в~тезисах докладов.
        \else%
            .
        \fi%
        \ifnum \value{citeregistered}=1%
            \ifnum \value{citeauthorpatent}=1%
                Зарегистрирован \arabic{citeauthorpatent} патент.
            \fi%
            \ifnum \value{citeauthorprogram}=1%
                Зарегистрирована \arabic{citeauthorprogram} программа для ЭВМ.
            \fi%
        \fi%
        \ifnum \value{citeregistered}>1%
            Зарегистрированы\ %
            \ifnum \value{citeauthorpatent}>0%
            \formbytotal{citeauthorpatent}{патент}{}{а}{}%
            \ifnum \value{citeauthorprogram}=0 . \else \ и~\fi%
            \fi%
            \ifnum \value{citeauthorprogram}>0%
            \formbytotal{citeauthorprogram}{программ}{а}{ы}{} для ЭВМ.
            \fi%
        \fi%
        % К публикациям, в которых излагаются основные научные результаты диссертации на соискание учёной
        % степени, в рецензируемых изданиях приравниваются патенты на изобретения, патенты (свидетельства) на
        % полезную модель, патенты на промышленный образец, патенты на селекционные достижения, свидетельства
        % на программу для электронных вычислительных машин, базу данных, топологию интегральных микросхем,
        % зарегистрированные в установленном порядке.(в ред. Постановления Правительства РФ от 21.04.2016 N 335)
    \end{refsection}%
    \begin{refsection}[bl-author, bl-registered]
        % Это refsection=2.
        % Процитированные здесь работы:
        %  * попадают в авторскую библиографию, при usefootcite==0 и стиле `\insertbiblioauthorimportant`.
        %  * ни на что не влияют в противном случае
        \nocite{vakbib2}%vak
        \nocite{patbib1}%patent
        \nocite{progbib1}%program
        \nocite{bib1}%other
        \nocite{confbib1}%conf
    \end{refsection}%
        %
        % Всё, что вне этих двух refsection, это refsection=0,
        %  * для диссертации - это нормальные ссылки, попадающие в обычную библиографию
        %  * для автореферата:
        %     * при usefootcite==0, ссылка корректно сработает только для источника из `external.bib`. Для своих работ --- напечатает "[0]" (и даже Warning не вылезет).
        %     * при usefootcite==1, ссылка сработает нормально. В авторской библиографии будут только процитированные в refsection=0 работы.
}

При использовании пакета \verb!biblatex! будут подсчитаны все работы, добавленные
в файл \verb!biblio/author.bib!. Для правильного подсчёта работ в~различных
системах цитирования требуется использовать поля:
\begin{itemize}
        \item \texttt{authorvak} если публикация индексирована ВАК,
        \item \texttt{authorscopus} если публикация индексирована Scopus,
        \item \texttt{authorwos} если публикация индексирована Web of Science,
        \item \texttt{authorconf} для докладов конференций,
        \item \texttt{authorpatent} для патентов,
        \item \texttt{authorprogram} для зарегистрированных программ для ЭВМ,
        \item \texttt{authorother} для других публикаций.
\end{itemize}
Для подсчёта используются счётчики:
\begin{itemize}
        \item \texttt{citeauthorvak} для работ, индексируемых ВАК,
        \item \texttt{citeauthorscopus} для работ, индексируемых Scopus,
        \item \texttt{citeauthorwos} для работ, индексируемых Web of Science,
        \item \texttt{citeauthorvakscopuswos} для работ, индексируемых одной из трёх баз,
        \item \texttt{citeauthorscopuswos} для работ, индексируемых Scopus или Web of~Science,
        \item \texttt{citeauthorconf} для докладов на конференциях,
        \item \texttt{citeauthorother} для остальных работ,
        \item \texttt{citeauthorpatent} для патентов,
        \item \texttt{citeauthorprogram} для зарегистрированных программ для ЭВМ,
        \item \texttt{citeauthor} для суммарного количества работ.
\end{itemize}
% Счётчик \texttt{citeexternal} используется для подсчёта процитированных публикаций;
% \texttt{citeregistered} "--- для подсчёта суммарного количества патентов и программ для ЭВМ.

Для добавления в список публикаций автора работ, которые не были процитированы в
автореферате, требуется их~перечислить с использованием команды \verb!\nocite! в
\verb!Synopsis/content.tex!.
